% ID: SEL-D0064-OA-E04
% ESTADO: cotejada. Segunda pasada visual y borrador compilado revisados.
% Fuente: Selectividad/sel_2010_matematicas_aplicadas/mat_apli/mat_apli_sept/mat_apli_sept/EXAMEN-ANDALUCÍA-3.pdf
% SHA-256: 9c3d2d628115a2f20f17b47171106ddd90f214b5f6804bdf9a7dcb7037c75efd
% Página: 1; opción A; ejercicio original 4.

\textbf{(2.5 puntos)} En una determinada especie animal el porcentaje de mortalidad debida a una enfermedad vírica es de al menos un $40\%$.

Se está realizando un estudio para probar la eficacia de un fármaco que permite tratar esa enfermedad y, consecuentemente, reducir el porcentaje de mortalidad en esa especie. Para ello, se suministró el fármaco a 50 sujetos enfermos, elegidos al azar, de los que murieron 14.

A la vista de estos datos, y tomando como nivel de significación $0.015$, \textquestiondown{}se puede afirmar que existe evidencia estadística suficiente para rechazar la hipótesis $H_0:p\geq0.4$, donde $p$ es la proporción, y por lo tanto aceptar la eficacia del fármaco?
